\chapter{Introduzione}
\label{cap:introduzione}

Il contesto applicativo di cui si tratta in questo lavoro è quello dell'\gls{iot}, un paradigma che si riferisce alla connessione di oggetti fisici a Internet, consentendo loro di comunicare e interagire tra loro e con gli utenti.
L'\gls{iot} ha un impatto significativo su diversi settori, tra cui l'agricoltura, dove può essere utilizzato per monitorare le condizioni del terreno, gestire l'irrigazione, monitorare la salute delle colture e ottimizzare l'uso delle risorse. 
In questo contesto, le tecnologie di comunicazione a bassa potenza e lunga distanza, come LoRaWAN, sono particolarmente adatte per la trasmissione dei dati raccolti dai sensori in ambienti rurali, dove la copertura di rete può essere limitata.

\section{L'azienda}
\textbf{M31} nasce nel 2007 ad opera di un gruppo di imprenditori, ricercatori universitari e professionisti appassionati di tecnologia, coordinati da Ruggero Frezza.
\\
Sin dalla sua costituzione, il cuore di M31 è il suo team di ingegneria - nei quindici anni di operatività, la Società ha svolto più di cento progetti di innovazione a vantaggio di startup proprie, imprese del territorio e grandi gruppi.
\\
La visione aziendale è quella di affiancare i clienti nel realizzare le loro idee di innovazione mediante l'applicazione di tecnologie allo stato dell'arte e degli output della ricerca universitaria.
\\
Tutto questo incarna il loro motto: \textbf{«With You. Beyond.»}
\\ \cite{site:sito-m31}.

\section{L'idea}
L'idea di questo lavoro è progettare, realizzare e validare sperimentalmente una infrastruttura LoRaWAN privata completa orientata a contesti agricoli. 
Il sistema, infatti, prevede l'implementazione di tutta l'infrastruttura:
\begin{itemize}
    \item dallo studio, dalla scelta e dall'implementazione dei nodi sensore LoRaWAN con firmware custom per la raccolta dati,
    \item passando per un gateway LoRaWAN implementato su Raspberry Pi con concentratore SX1303 in banda EU868;
    \item fino ad arrivare al server LoRaWAN locale basato su ChirpStack, per la gestione dei dispositivi, la raccolta dei dati e l'integrazione con sistemi esterni.
\end{itemize}
Il lavoro non si limita alla sola installazione dei componenti, ma punta a caratterizzare il sistema dal punto di vista prestazionale, individuandone limiti, comportamento reale in campo e linee guida per una possibile industrializzazione.

\section{Organizzazione del testo}
% TODO: aggiornare i titoli dei capitoli
\begin{description}
    \item[{\hyperref[cap:processi-metodologie]{Il secondo capitolo}}] descrive ...
    
    \item[{\hyperref[cap:descrizione-stage]{Il terzo capitolo}}] approfondisce ...
    
    \item[{\hyperref[cap:analisi-requisiti]{Il quarto capitolo}}] approfondisce ...
    
    \item[{\hyperref[cap:progettazione-codifica]{Il quinto capitolo}}] approfondisce ...
    
    \item[{\hyperref[cap:verifica-validazione]{Il sesto capitolo}}] approfondisce ...
    
    \item[{\hyperref[cap:conclusioni]{Nel settimo capitolo}}] descrive ...
\end{description}

% TODO: aggiornare le convenzioni
Riguardo la stesura del testo, relativamente al documento sono state adottate le seguenti convenzioni tipografiche:
\begin{itemize}
	\item gli acronimi, le abbreviazioni e i termini ambigui o di uso non comune menzionati vengono definiti nel glossario, situato alla fine del presente documento;
	\item per la prima occorrenza dei termini riportati nel glossario viene utilizzata la seguente nomenclatura: \emph{parola}\glsfirstoccur;
	\item i termini in lingua straniera o facenti parti del gergo tecnico sono evidenziati con il carattere \emph{corsivo}.
\end{itemize}
